% Level 1, Answers week 12 -- Showcase and Celebration   (input by answer-key.tex)
\answersweek{12}{Showcase and Celebration}{Notes for workbook chapter 12 and homework 12. Exit assessment answers are in the Assessments section of this key.}

\section{Workbook}
\begin{answer}[Exercise 12.1: Prepare your 2-minute demo]
A good script names the project, lets the visitor answer, connects one answer to the \texttt{if} (\enquote{because you said night, it drew a moon}), points at the function and the loop, and mentions one real bug. Use the \enquote{Explaining} row of the rubric in the Project Teacher Guide.
\end{answer}

\begin{answer}[Exercise 12.2: Be a great visitor]
Any idea they admired, said specifically (\enquote{Mei's fish had random sizes so it looked like a real sea}).
\end{answer}

\begin{answer}[Exercise 12.3: What I can do now]
The eight sentences match the Level 1 goals in order: L1.G1 algorithms, L1.G2 loops, L1.G3 conditionals, L1.G4 variables/input/random, L1.G5 functions, L1.G6 debugging, L1.G7 pair work and explaining, L1.G8 safety and credit. Unticked sentences are a useful note for the start of Level 2 --- not a failure.
\end{answer}

\section{Homework}
\begin{answer}[Core: My Level 1 reflection]
Open answers. Look for specifics (\enquote{I can read error messages now}) and a wish for Level 2 you can use when planning (games, art, quizzes).
\end{answer}

\begin{answer}[Challenge: holiday menu (sample solutions)]
Menu 1 --- Magic spiral:
\begin{pycode}
import turtle

t = turtle.Turtle()
t.speed(0)
t.color("#0072B2")

for i in range(60):
    t.forward(i * 3)
    t.left(91)

turtle.done()
\end{pycode}
Menu 3 --- Dice roller:
\begin{pycode}
import random

print("Dice roller!")
roll = random.randint(1, 6)
print("You rolled", roll)
if roll == 6:
    print("A six! Roll again!")
\end{pycode}
Menu 4 --- City skyline (the drawing part):
\begin{pycode}
def building(height, color):
    t.pendown()
    t.color("#000000", color)
    t.begin_fill()
    for side in range(2):
        t.forward(50)
        t.left(90)
        t.forward(height)
        t.left(90)
    t.end_fill()
    t.penup()
    t.forward(60)


for i in range(12):
    building(random.randint(50, 300), random.choice(PALETTE))
\end{pycode}
(Before the loop: \texttt{import random}, the \texttt{PALETTE} list, a turtle \texttt{t} and \texttt{t.penup()}, \texttt{t.goto(-350, -200)}; after it: \texttt{turtle.done()}.) Menu 2 is the Square Flower with \texttt{range(36)}, \texttt{t.left(10)} and \texttt{random.choice(PALETTE)}. Menu 5 extends homework 9 with more questions.
\end{answer}

\begin{answer}[Creative: invent a challenge]
Any clear challenge. Bonus if it names the ingredients it needs (\enquote{a loop and random}) --- that is decomposition.
\end{answer}
